\section{Immutable MVP、Gossip 与 Day 1/100/1000}

Vercel 的早期 MVP 极窄：一个 API/CLI 调用返回不可变 deployment URL。底层系统起初可能过度配置且昂贵，但这个接口足以验证用户是否真正想要“部署即 URL”。Rauch 将 immutable data structures、IPFS content identity 和 gossip propagation 作为灵感来源，随后用实际需求驱动内部重构。

\subsection{先验证 idea-market fit，再优化内部实现}

\term{immutable deployment} 指部署内容一旦生成便不原地修改；新版本创建新 identity。\term{gossip protocol} 让节点互相传播 metadata，最终在全球副本间收敛。Domain 指针可以变化，deployment artifact 保持不变，从而获得可追溯 preview、原子 promote 和快速 rollback。

\lecturefigure{15-immutable-mvp.png}{一个窄 API/CLI 把 artifact identity、metadata propagation 和 URL 连接起来。}{本地字幕 32:50--35:55；概念重绘。}

\begin{knowledgebox}{读图：MVP 可以内部低效，但外部契约要正确}
早期系统可以 over-provision，只要成本可承受且产品信号值得验证；但若外部 API 把可变服务器状态暴露给用户，后续很难迁移。优秀 MVP 选择一个未来仍成立的稳定抽象，再逐步替换内部实现。
\end{knowledgebox}

\subsection{Day 1、Day 100、Day 1000 是三种不同验收}

上一小节说明 MVP 可以暂时内部低效，但外部契约必须长期正确；接下来需要用时间尺度检验这个契约是否真的吸收了复杂度。Day 1 检查 onboarding 和第一印象；Day 100 检查项目增长后平台是否泄漏复杂性；Day 1000 检查 uptime、恢复、成本和组织信任。许多开发工具只优化 Day 1 demo，却把长期状态、依赖和运维负担交还给用户，因此成熟度不能只由首次部署速度衡量。

\lecturefigure{16-day-maturity.png}{基础设施产品必须分别通过采用、复杂度吸收和长期运营测试。}{本地字幕 36:00--37:05；概念重绘。}

\begin{importantbox}{Day 1/100/1000 验收问题}
\begin{itemize}
  \item Day 1：新人能否在不了解底层细节时安全部署？
  \item Day 100：route、team、environment 和数据依赖增加后，默认模型是否仍清晰？
  \item Day 1000：升级、事故、人员流动和供应商变化时，系统能否恢复并保持可审计？
\end{itemize}
\end{importantbox}

\begin{warningbox}{不可变部署不等于不可变数据}
代码和 assets 可以 immutable，数据库、队列和外部系统仍会变化。Rollback application version 可能遇到 schema incompatibility。平台需要 migration policy、backward-compatible contract 和 data recovery，不能把 Git 式体验机械套在所有状态上。
\end{warningbox}

\subsection{本章小结}

好的 MVP 固定正确的外部抽象，允许内部暂时低效；好的基础设施产品则继续通过 Day 100 和 Day 1000，吸收项目复杂度并保持运营信任。

